\bookchapter{第四章：アプリのインストールと管理——スマートフォンの「スキルパック」}

\chapterintro{アプリがスマートフォンを万能にするが、管理も必要だ。新入社員に席を用意し権限を割り当てるように、アプリのインストールにも規範フローが必要だ。}

\booksection{一、公式ストア：最も安全な「アプリスーパー」}

\begin{itemize}
  \item \textbf{Android}iOS
  \item \textbf{iOS}：App Storeのみ。
  \item \textbf{HarmonyOS}コアのセキュリティ意義：審査済み、更新保証あり、決済安全。行動ガイド：
\end{itemize}

公式ストアからのダウンロードを優先！\textbf{二、ピュア/プロテクトモード：スマートフォンの「アンチウイルス壁」を有効化}

\booksection{二、ピュア/プロテクトモード：スマートフォンの「アンチウイルス壁」を有効化二、ピュア/プロテクトモード：スマートフォンの「アンチウイルス壁」を有効化}

ピュアモード\textbf{または}プロテクトモード\textbf{を内蔵し、これはシステムレベルのセキュリティスイッチだ。}すべてのユーザーに有効化を強く推奨\textbf{。}。

\booksubsection{それは一体何？}

コア機能はサードパーティアプリのインストールをブロック\textbf{すること：有効時、システムは}公式アプリストアからのアプリインストールのみ許可\textbf{し、ブラウザ、QRコード、チャットソフトなどのチャネルからダウンロードしたインストーラは直接ブロックされ、悪意あるソフトを根源から排除する。}2. 適用対象

\booksubsection{適用対象}

\begin{itemize}
  \item \textbf{：「公式」と「偽物」の広告を区別できず、ピュアモードが最も確実なお守り。}アプリ権限に敏感でないユーザー
  \item \textbf{：ニッチアプリ、ベータ版、海外版ゲームをインストールする必要がない。}デバイスのセキュリティを向上させたいすべてのユーザー
  \item \textbf{：マニアであっても、有効にしておく方が無効より良い——一時的な必要があれば、その時だけ無効にすればよい。}3. 有効化パス（Huawei/Xiaomiの場合）
\end{itemize}

\booksubsection{有効化パス（Huawei/Xiaomiの場合）}

\begin{itemize}
  \item \textbf{：[設定] > [システムと更新] >}[ピュアモード]\textbf{、「有効化」を選択。}、「オンにする」を選択。
  \item \textbf{Xiaomi}[プライバシー保護]\textbf{または [設定] > [アプリ設定] > [アプリインストール]、「提供元不明アプリのインストール」をオフ。}OPPO/vivo
  \item \textbf{：一部機種はアンロック可能、申請が必要。}[外部提供元アプリのインストール]\textbf{、個別に認証をオフ。}4. アプリストアのセキュリティ検査メカニズム
\end{itemize}

\booksubsection{アプリストアのセキュリティ検査メカニズム}

ウイルススキャン

\begin{itemize}
  \item \textbf{：各アプリは上架前に複数のアンチウイルスエンジンによる合同スキャンを受け、既知のトロイの木馬やウイルスの特徴を識別。}権限審査
  \item \textbf{：ストアはアプリが申請する権限が「合理的」かを審査——懐中電灯アプリが連絡先権限を要求する？即座に却下または下架。}開発者認証
  \item \textbf{：上架には企業資格、個人身分証明またはクレジットカード認証が必要。問題発生時に人物を特定でき、「ヤミアプリ」ではない。}継続監視
  \item \textbf{：上架済みアプリも継続的に再スキャンされ、新たに発見された悪意ある挙動は強制下架される。}三、サードパーティ提供元と「電子寄生虫」：リスク極めて高い！
\end{itemize}

\booksection{三、サードパーティ提供元と「電子寄生虫」：リスク極めて高い！三、サードパーティ提供元と「電子寄生虫」：リスク極めて高い！}

「電子寄生虫」\textbf{。}。

\textbf{トロイの木馬が植入された「クラッキング版」ソフトウェア}\par

を指す——あなたは「節約した」と思うが、実際は「罠にはまった」。一般的な形式：\textbf{「クラッキング版」「無料版」「広告除去版」の有料ソフトウェア；}「いいね増やし」「紅包搶奪」「ゲームチート」などのグレーツール；

\begin{itemize}
  \item グループチャットの「広告」が推奨する「内部版」アプリ。
  \item これらのプログラムは表面上は使えるが、裏で
  \item 連絡先を盗み、SMSを読み、マイクを盗聴、バックグラウンドでマイニング
\end{itemize}

している可能性がある。ピュアモードはそれらを防ぐ\textbf{第一の防線}——そもそもインストールさせない。\textbf{第二の防線の推奨：}定期全スキャン

\booksubsection{第二の防線の推奨：}

\begin{itemize}
  \item \textbf{センシティブ権限の慎重な付与}：連絡先、SMS、位置、カメラ、マイク——これら5項目の「高危権限」は厳格に審査必要。
  \item \textbf{消費電力とデータ通信量の異常への注目}：スマートフォンが突然発熱、バッテリー急消費、深夜のデータ通信——八割バックグラウンドに「寄生虫」が隠れている。
  \item \textbf{万が一サードパーティ製アプリをインストールする必要がある場合：}提供元が絶対に信頼できることを確認する。
\end{itemize}

インストール前に権限要求を注意深く確認する。

\begin{enumerate}
  \item インストール後にセキュリティアプリでスキャンする。
  \item 最終的なアドバイス：システム設定で「提供元不明のアプリのインストール」権限をオフにする。
  \item 四、インストール、更新、アンインストール
\end{enumerate}

インストール

\booksection{四、インストール、更新、アンインストール四、インストール、更新、アンインストール}

\begin{itemize}
  \item \textbf{更新}：
  \item \textbf{必ず更新すること！}：\textbf{アンインストール}：アイコンを長押ししてアンインストール領域にドラッグ。より完全に削除したい場合は、[設定] > [アプリ管理]で該当アプリを見つけてアンインストール。
  \item \textbf{五、権限管理：あなたはアプリの「ボス」}アプリが権限（連絡先や位置など）を要求するのは、合理的な場合（地図アプリが位置を要求）もあれば、不合理な場合（懐中電灯が連絡先を要求）もある。
\end{itemize}

\booksection{五、権限管理：あなたはアプリの「ボス」五、権限管理：あなたはアプリの「ボス」}

：「常に許可」「使用中のみ許可」「禁止」「毎回確認」。

\begin{itemize}
  \item \textbf{管理原則}：
  \item \textbf{最小化：不要なものは与えない。}：
  \begin{enumerate}
    \item 状況に応じた判断：タクシー配車アプリのように、使用時のみ位置権限を与える。
    \item 重点監視すべき権限
    \item 連絡先/SMS：迷惑SMS送信に悪用される可能性。
  \end{enumerate}
\end{itemize}

\begin{tcbwarn}{リアルタイム位置：行動追跡される可能性。}
\begin{itemize}
  \item カメラ/マイク：盗撮·盗聴される可能性。
  \item ファイルアクセス：プライベートファイルが読み取られる可能性。
  \item 六、ソフトウェアインストールパッケージの構成要素
  \item アプリストアからダウンロードするインストールパッケージの中に何が入っているか考えたことがあるだろうか？インストールパッケージの構成を理解することで、アプリの仕組みをより良く理解でき、問題発生時にどこから調査すべきかも分かるようになる。
\end{itemize}

\end{tcbwarn}

\booksection{Android APKファイル構造}

APK（Android Package Kit）は本質的にZIP圧縮ファイルで、解凍すると以下のコアディレクトリとファイルが見える：

\booksubsection{Android APKファイル構造}

：アプリの「身分証明書」であり「取扱説明書」。アプリのパッケージ名、バージョン番号、必要な権限、サポートするシステムバージョン、入口Activityなどのコア情報を記録。システムがアプリをインストールする際、最初にこのファイルを読み取る。

\begin{itemize}
  \item \textbf{classes.dex}：アプリのコアコード。Java/KotlinソースコードがコンパイルされてDalvikバイトコードになり、このファイルに格納される。アプリによっては複数のdexファイル（classes2.dex、classes3.dexなど）を持つ場合がある。
  \item \textbf{resources.arsc}：コンパイル済みのリソースインデックステーブル。すべてのリソースのIDと対応関係を記録。
  \item \textbf{res/ディレクトリ}：リソースフォルダ。画像、レイアウトファイル、文字列、色定義などを含む。異なる解像度の画像は異なるdrawableフォルダ（drawable-hdpi、drawable-xhdpiなど）に配置される。
  \item \textbf{assets/ディレクトリ}：生リソースフォルダ。コンパイル不要のリソース（ウェブファイル、フォント、オーディオなど）を格納。アプリ実行時にAssetManagerで読み取り可能。
  \item \textbf{lib/ディレクトリ}：ネイティブライブラリフォルダ。コンパイル済みの.soファイル（C/C++コードのコンパイルの産物）をCPUアーキテクチャ別（arm64-v8a、armeabi-v7aなど）に格納。ゲームや高性能計算を必要とするアプリには通常このディレクトリがある。
  \item \textbf{META-INF/ディレクトリ}：署名情報フォルダ。MANIFEST.MF、CERT.SF、CERT.RSAの3ファイルを含み、APKの完全性と署名の検証に使用。
  \item \textbf{2. iOS IPAファイル構造}IPA（iOS App Store Package）も同様に圧縮ファイルだが、構造は若干異なる：
\end{itemize}

\booksubsection{iOS IPAファイル構造}

：コアコンテンツの格納場所。内部に.appフォルダがあり、アプリのすべてのコンテンツを含む。

\begin{itemize}
  \item \textbf{Info.plist}：アプリの設定ファイルで、AndroidのAndroidManifest.xmlに類似。アプリID、バージョン、権限要件、サポートデバイスなどを記録。
  \item \textbf{実行可能ファイル}：.appフォルダと同名のメインプログラムで、コンパイル済みのMach-O形式バイナリファイル。
  \item \textbf{Frameworks/ディレクトリ}：アプリが依存する動的ライブラリとフレームワーク。
  \item \textbf{リソースファイル}：画像、オーディオ、動画などのリソース。通常.appフォルダ内に格納。
  \item \textbf{\_CodeSignature/ディレクトリ}：コード署名情報。アプリの完全性と出所の検証に使用。
  \item \textbf{3. 署名検証メカニズム}なぜ公式ストアからダウンロードするアプリの方が安全なのか？署名検証が鍵だ。
\end{itemize}

\booksubsection{署名検証メカニズム}

：開発者が秘密鍵でAPKに署名し、システムインストール時に公開鍵で検証する。APKが改ざんされた場合（ウイルスの植入など）、署名は無効になり、システムはインストールを拒否する。Google PlayはさらにApp Signingを要求し、Googleが署名鍵を保管することでアプリのセキュリティをさらに保護する。

\begin{itemize}
  \item \textbf{iOS署名}：Appleの審査はより厳格。各IPAにはAppleが発行した証明書の署名が必須で、なければインストール不可。これがiOSがApp Storeまたは企業証明書からしかアプリをインストールできない理由だ。企業証明書で署名されたアプリは有効期限があり、Appleはいつでも取り消し可能。
  \item \textbf{署名の役割}：
  \item \textbf{出所の検証：アプリが自称する開発者から実際に来ていることを確認。}：
  \begin{enumerate}
    \item なりすまし防止：悪意あるアプリが他のアプリに偽装することを防ぐ。
    \item 4. インストールパッケージのサイズと最適化
    \item アプリストアに表示されるアプリサイズと、インストール後に占有する容量が異なることに気づいたことがあるかもしれない。理由は以下の通り：
  \end{enumerate}
\end{itemize}

\booksubsection{インストールパッケージのサイズと最適化}

：圧縮後のAPK/IPAのサイズ。

\begin{itemize}
  \item \textbf{インストールサイズ}：解凍後のサイズで、通常ダウンロードサイズよりはるかに大きい。
  \item \textbf{実行時占有容量}：アプリ実行時にキャッシュやデータファイルが生成され、さらに容量を占有する。
  \item \textbf{現代アプリの最適化手法：}App Bundle（Android）
\end{itemize}

：Google Playの動的配信技術。ユーザーはダウンロード時に自分のデバイスに適したリソースのみを取得（画面解像度に対応する画像のみなど）し、ダウンロードサイズを大幅に削減。

\begin{itemize}
  \item \textbf{App Thinning（iOS）}：類似の技術で、Appleがデバイスに応じて不要なリソースを自動的に切り詰める。
  \item \textbf{オンデマンドダウンロード}：ゲームなどの大型アプリは、まずコアパッケージをダウンロードし、後からリソースパッケージをダウンロード可能。
  \item \textbf{豆知識：なぜあるアプリのインストールパッケージは特に大きいのか？}ゲーム系アプリは通常最も大きく、大量のHDテクスチャ、3Dモデル、オーディオファイルを含むため。ソーシャル系アプリも比較的大きく、各種機能モジュールをサポートする必要があるため。ツール系アプリは通常最も小さい。あるツール系アプリが特に大きい場合、不要な機能がバンドルされていないか警戒すべき。
\end{itemize}

\begin{tcbinfo}{実践セクション：3つのアプリの権限を審査}
よく使うアプリを3つ（ソーシャル、ショッピング、ツール系など）選び、[設定]のアプリ権限管理ページで項目ごとに審査し、不要なセンシティブ権限をオフにするか「使用中のみ許可」に変更する。

\end{tcbinfo}

\begin{tcbpractice}{第五章：スマートフォンをもっと理解させる——基本設定}
数分かけて設定するだけで、体験は大幅に向上する。これらの設定はスマートフォンに「チューニング」を行うようなもので、あなたの使用習慣に合わせることができる。

\end{tcbpractice}
